% Q13 -- Fault models: (a) stuck-at fault on a gate output (S-A-0 / S-A-1),
% (b) bridging short between two adjacent nets (S1, S2), (c) open in an nMOS source-to-GND.
\documentclass[border=10pt]{standalone}
\usepackage{circuitikz}
\begin{document}
\begin{circuitikz}[font=\small]
% ---------- (a) Stuck-at fault ----------
\node[nand port, scale=1.1] (g) at (0,5) {};
\draw (g.in 1) -- ++(-0.5,0) node[left]{$A$};
\draw (g.in 2) -- ++(-0.5,0) node[left]{$B$};
\draw (g.out) -- ++(0.8,0) coordinate (z);
\node[circ] at (z){};
\node[right, align=left, font=\scriptsize] at ([xshift=2pt]z) {output node\\ Stuck-At-0 / Stuck-At-1};
\node at (0.2,3.3) {(a) Stuck-at fault};
% ---------- (b) Bridging / short fault ----------
\draw (5,5.6) -- (8,5.6) node[right]{net A ($S_1\Rightarrow$ S-A-0)};
\draw (5,4.4) -- (8,4.4) node[right]{net B};
\draw[red, thick] (6.2,5.6) -- (6.2,4.4);
\node[circ] at (6.2,5.6){}; \node[circ] at (6.2,4.4){};
\node[red, above right, font=\scriptsize] at (6.25,5.0) {bridging short $S_1$/$S_2$};
\node at (6.5,3.3) {(b) Bridging / short fault};
% ---------- (c) Open fault ----------
\path (12,5.2) node[nmos] (m){};
\draw (m.D) -- (12,6.4) node[above]{out};
\draw (m.G) -- (10.8,5.2) node[left]{in};
\draw (m.S) -- (12,4.3) coordinate (s1);
% broken (open) source-to-GND connection
\draw (12,3.7) -- (12,3.1) node[ground]{};
\draw[red, thick] (11.7,4.15) -- (12.3,3.85);
\node[red, right, font=\scriptsize] at (12.25,4.0) {Open};
\node at (12,2.5) {(c) Open fault};
\end{circuitikz}
\end{document}
